\bookchapter{الفصل 62: تطوير الواجهات — مدخل إلى SwiftUI}

\chapterintro{الواجهة هي كل ما يراه المستخدم. SwiftUI إطار واجهات «تقريري» أطلقته Apple عام 2019 — فأنت تحتاج فقط إلى «وصف» شكل الواجهة، ويتولى النظام تلقائيًا الرسم والتكيف مع الوضع الداكن والشاشات المختلفة. بحلول عام 2026 أصبح SwiftUI السائد بلا منازع في المشاريع الجديدة. في هذا الفصل سنتعلم «الرسم».}

\booksection{أولًا: فكرة SwiftUI — تقريري لا أمري}

تطوير الواجهات القديم «أمري»: تخبر النظام خطوة خطوة «أنشئ زرًا ضعه عند (20,30) اضبط اللون». أما SwiftUI «تقريري»: تصف مباشرة «في الواجهة زر أحمر مكتوب عليه ‘اضغطني’»، أما كيف يُرسم ومتى يُعاد رسمه فالنظام يتكفل بذلك بنفسه.

\begin{lstlisting}[language=Swift]
import SwiftUI

struct ContentView: View {
    var body: some View {      // body: يصف "كيف تبدو الواجهة"
        VStack {               // مرتبة رأسيًا
            Text("مرحبًا يا عالم!")
                .font(.largeTitle)      // معدِّل: عنوان كبير
            Button("اضغط هنا") {
                print("تم الضغط!")
            }
            .buttonStyle(.borderedProminent)  // زر أزرق ممتلئ
        }
        .padding()
    }
}
\end{lstlisting}

النقاط الأساسية: العرض (View) بنية تلتزم ببروتوكول View، والخاصية body تصف الواجهة، والمُعدِّلات على الصيغة .modifier() تتسلسل «كالمكياج» طبقة فوق طبقة. هذا هو كامل أساس SwiftUI.

\booksection{ثانيًا: العروض الأساسية — نص وزر وحقل إدخال وصورة}

\begin{center}
\begin{tabularx}{\linewidth}{|X|X|X|}
\hline
\textbf{العرض} & \textbf{وظيفته} & \textbf{مثال} \tabularnewline
\hline
Text & عرض نص & Text("مرحبًا") \tabularnewline
Button & الزر & Button("اضغطني") \{ الإجراء \} \tabularnewline
TextField & حقل الإدخال & TextField("تلميح", text: \$النص) \tabularnewline
Image & عرض صورة & \texttt{Image(systemName: "star.\allowbreak fill")} \tabularnewline
\hline
\end{tabularx}
\end{center}

\begin{center}\small\color{gray}العروض الأساسية الأربعة في SwiftUI\end{center}

لاحظ Image(systemName:) — فقد زودت Apple مكتبة أيقونات مدمجة اسمها SF Symbols (آلاف الأيقونات المجانية)، واسم مثل “star.fill” يعرض مباشرة نجمة مصمتة. هذه ميزة خفية في تطوير iOS: تكاد لا تحتاج إلى البحث عن ملفات أيقونات عند بناء الواجهات.

\booksection{ثالثًا: ثالوث التخطيط — VStack وHStack وZStack}

\begin{itemize}
  \item VStack: ترتيب عمودي (كالاصطفاف في طابور الطعام)؛
  \item HStack: ترتيب أفقي (كالجلوس جنبًا إلى جنب)؛
  \item ZStack: تراكب (كأوراق اللعب المكدسة، مناسب للشارات والطبقات العلوية).
\end{itemize}

\begin{lstlisting}[language=Swift]
VStack(alignment: .leading, spacing: 12) {
    HStack {
        Text("الحساب")
        Spacer()                 // مساحة مرنة: تدفع ما على اليمين إلى النهاية
        Text("-25.5 يوان").foregroundColor(.red)
    }
    Divider()                    // خط فاصل
    Text("ملاحظة: الغداء")
}
.padding(16)                     // هامش من كل الجهات
.background(Color(.systemGray6)) // خلفية رمادية فاتحة
.cornerRadius(12)                // زوايا مستديرة

// ترتيب المعدلات مهم: padding أولًا ثم background؛ الخلفية لا تشمل الهامش
\end{lstlisting}

\begin{tcbwarn}{تنبيه: ترتيب المُعدِّلات له شأن}
أثر padding(16).background(رمادي) يختلف عن background(رمادي).padding(16) — الأولى «تترك فراغًا ثم تلوّن» (الخلفية لا تشمل الهوامش)، والثانية «تلوّن ثم تترك فراغًا» (الخلفية تشمل الهوامش). ومن هنا يأتي «الحاشية البيضاء الغريبة» التي تربك المبتدئين أكثر من أي شيء آخر. وإذا شككت فبدّل الترتيب، وستجيبك المعاينة فورًا.

\end{tcbwarn}

\booksection{رابعًا: الحالة والربط — كيف تنبض الواجهة بالحياة}

الواجهة الثابتة بلا طعم. قاعدة SwiftUI: الواجهة تعكس الحالة بأمانة دائمًا، وإذا تغيرت الحالة أعادت الواجهة رسم نفسها تلقائيًا.

\begin{lstlisting}[language=Swift]
struct CounterView: View {
    @State private var count = 0   // @State: إذا تغيرت هذه القيمة يجب إعادة الرسم

    var body: some View {
        VStack {
            Text("ضغطت \(count) مرات")
            Button("اضغط +1") {
                count += 1          // غيّر الحالة وستتحدث الواجهة تلقائيًا
            }
        }
    }
}
\end{lstlisting}

حقل الإدخال يحتاج «ربطًا ثنائي الاتجاه» — فالرمز \$ يمرر «القيمة وحق تعديلها» معًا:

\begin{lstlisting}[language=Swift]
struct InputView: View {
    @State private var text = ""

    var body: some View {
        TextField("اكتب شيئًا", text: $text)  // ربط ثنائي الاتجاه بـ $
            .textFieldStyle(.roundedBorder)
        Text("كتبت: \(text)")   // يتحدث فوريًا أثناء الكتابة
    }
}
\end{lstlisting}

عملية المستخدم

الحالة تتغير (count += 1)

SwiftUI يلاحظ التغير

يعيد رسم الواجهة (إعادة حساب body)

\begin{center}\small\color{gray}شكل 3: تدفق البيانات أحادي الاتجاه في SwiftUI من «الحالة الواجهة»\end{center}

\booksection{خامسًا: المعاينة والسمات والحركة — «قوى SwiftUI الخارقة»}

\begin{itemize}
  \item ماكرو \#Preview: أضف \#Preview \{ ContentView() \} في أسفل الملف، فتعرض لوحة الجانب النتيجة لحظيًا — عدّل الكود وترى النتيجة فورًا؛
  \item الوضع الداكن مجانًا: استخدم ألوان النظام (Color(.systemGray6) و.primary) بدلًا من تثبيت الألوان يدويًا، فيتكيف الوضع الداكن تلقائيًا دون سطر إضافي واحد؛
  \item الحركة بسطر واحد: withAnimation \{ count += 1 \} أو أضف إلى العرض .animation(.default, value: count) — فتصبح الانتقالات انسيابية.
\end{itemize}

\begin{lstlisting}[language=Swift]
struct AnimatedHeartView: View {
    @State private var liked = false

    var body: some View {
        Button {
            withAnimation(.spring(duration: 0.4)) {
                liked.toggle()
            }
        } label: {
            Image(systemName: liked ? "heart.fill" : "heart")
                .font(.system(size: 48))
                .foregroundColor(liked ? .red : .gray)
                .scaleEffect(liked ? 1.3 : 1.0)
        }
    }
}
\end{lstlisting}

\begin{tcbinfo}{نصيحة: الوظيفة أولًا ثم الحركة}
من أكثر علات المبتدئين «الانشغال بالحركة قبل أن تصح الواجهة». اجعل الوظائف تعمل أولًا، ثم أضف المؤثرات الحركية دفعة واحدة في النهاية — فالحركة زينة تُضاف لما أُحسن، والوظيفة إغاثة تُقدَّم في وقت الحاجة. وهذا هو مبدؤنا الدائم: البساطة فوق كل شيء.

\end{tcbinfo}

\begin{tcbpractice}{هيا نجرّب: اصنع «زر المزاج»}
\booksection{سادسًا: مرجع سريع للمُعدِّلات الشائعة — صفحة بيدك فلا واجهة تقلقك}

\begin{center}
\begin{tabularx}{\linewidth}{|X|X|}
\hline
\textbf{ماذا تريد أن تفعل} & \textbf{المُعدِّل} \tabularnewline
\hline
التباعد الداخلي & \texttt{.\allowbreak padding(16)} \tabularnewline
الزوايا الدائرية & \texttt{.\allowbreak cornerRadius(12)} \tabularnewline
لون الخلفية & \texttt{.\allowbreak background(Color(.\allowbreak systemGray6))} \tabularnewline
الخط وحجمه & \texttt{.\allowbreak font(.\allowbreak title.\allowbreak bold())}/\texttt{.\allowbreak font(.\allowbreak system(size: 20))} \tabularnewline
لون المقدمة (النص/الأيقونة) & \texttt{.\allowbreak foregroundColor(.\allowbreak red)} \tabularnewline
ملء العرض & \texttt{.\allowbreak frame(maxWidth: .\allowbreak infinity)} \tabularnewline
حجم ثابت & \texttt{.\allowbreak frame(width: 100, height: 44)} \tabularnewline
المحاذاة & \texttt{VStack(alignment: .\allowbreak leading)} \tabularnewline
الشفافية & \texttt{.\allowbreak opacity(0.\allowbreak 6)} \tabularnewline
الظل & \texttt{.\allowbreak shadow(radius: 5)} \tabularnewline
إيماءة النقر & \texttt{.\allowbreak onTapGesture \{ .\allowbreak .\allowbreak .\allowbreak  \}} \tabularnewline
نوع لوحة المفاتيح & \texttt{.\allowbreak keyboardType(.\allowbreak decimalPad)} \tabularnewline
\hline
\end{tabularx}
\end{center}

\begin{center}\small\color{gray}مرجع سريع للمُعدِّلات الشائعة في SwiftUI\end{center}

لا بأس إن لم تحفظها — فهذه الصفحة موجودة للرجوع إليها. الأكثر استخدامًا في الحقيقة هي الأربعة الأولى: padding والزوايا الدائرية والخلفية والخط.

حوّل المشروع النموذجي إلى مبدّل مزاج: زر واحد، وعند كل نقرة يتنقل بالتناوب بين سعيد حزين غاضب، ويتغير معه النص ولون الخلفية. تلميح: خزّن المجموعات الثلاث (نص + لون) في مصفوفة، وخزّن الفهرس في @State، وعند النقر استخدم (index + 1) \% 3. بعد الانتهاء أرسله لأصدقائك يلعبون به، واجمع موجة «واو» الأولى: «واو، أنت تصنع تطبيقات!».

\end{tcbpractice}

\begin{tcbsummary}{خلاصة الفصل}
\begin{itemize}
  \item SwiftUI تقريري: تصف «الشكل المطلوب»، والنظام يتولى الرسم وإعادة الرسم
  \item العرض بنية تلتزم ببروتوكول View؛ وbody يصف الواجهة؛ والمُعدِّلات تتسلسل «كالمكياج»
  \item العروض الأساسية الأربعة: Text / Button / TextField / Image؛ ومكتبة أيقونات SF Symbols مجانية بلا حدود
  \item ثالوث التخطيط: VStack عمودي وHStack أفقي وZStack متراكب؛ وSpacer مساحة مرنة
  \item @State تجعل الواجهة حية؛ والرمز \$ ربط ثنائي الاتجاه يوصل القيمة إلى حقل الإدخال
  \item \#Preview ما تراه هو ما تحصل عليه؛ وألوان النظام تتكيف تلقائيًا مع الوضع الداكن؛ وwithAnimation حركة بسطر واحد
\end{itemize}

\end{tcbsummary}
